\chapter{GitLab CLI: o GitLab pelo terminal}
\label{cap:glab}

\begin{conceito}[GitLab CLI]
O GitLab CLI, usado pelo comando \cmd{glab}, é a ferramenta oficial para \textbf{interagir com o GitLab diretamente pelo terminal}: Merge Requests, Issues, projetos, pipelines, branches, releases e autenticação.
\end{conceito}

\begin{center}
\begin{tikzpicture}[b/.style={caixa, minimum width=3.6cm, minimum height=1.4cm}]
  \node[b] (g) at (0,0) {\textbf{Git}\\\scriptsize controle de versão};
  \node[b=verdeMB] (l) at (4.6,0) {\textbf{GitLab}\\\scriptsize plataforma de colaboração};
  \node[b=ouroMB!85!black] (c) at (9.2,0) {\textbf{glab}\\\scriptsize o GitLab no terminal};
\end{tikzpicture}
\end{center}

\section{Instalação}

As formas de instalar mudam conforme o sistema e a versão. Siga a página oficial: \url{https://gitlab.com/gitlab-org/cli}, seção \emph{Installation}. Para conferir:

\begin{terminal}
glab --version
\end{terminal}

\section{Autenticação}

O \cmd{glab} precisa saber em qual GitLab e com qual usuário trabalhar. Num GitLab instalado pela instituição, informe o endereço dele:

\begin{terminal}
glab auth login --hostname <endereco-do-gitlab>
\end{terminal}

O comando pergunta como você quer se autenticar. O caminho mais comum é um \termo{token de acesso pessoal}, gerado no GitLab em \textbf{Preferências do usuário \faChevronRight\ Access tokens}.

\begin{atencao}[Token é senha]
Um token de acesso dá acesso à sua conta. Não compartilhe, não envie por e-mail ou chat e não grave em arquivo de projeto. Dê a ele só as permissões necessárias e uma data de validade.
\end{atencao}

\section{Os comandos mais úteis}

\begin{center}
\renewcommand{\arraystretch}{1.3}
\begin{tabularx}{\linewidth}{@{}>{\ttfamily\color{azulMB}}l X@{}}
\toprule
glab mr create & cria um Merge Request a partir da branch atual \\
glab mr create -s origem -b destino & escolhe as branches de origem e de destino \\
glab mr list & lista os Merge Requests abertos do projeto \\
glab mr view & mostra o Merge Request da branch atual \\
glab mr diff & mostra as alterações do Merge Request \\
glab mr approve & aprova um Merge Request (se você tiver permissão) \\
glab mr merge & faz o merge (se você tiver permissão) \\
glab issue list & lista as Issues \\
glab ci status & mostra a situação do pipeline da branch atual \\
glab repo clone grupo/projeto & clona um projeto pelo nome \\
\bottomrule
\end{tabularx}
\end{center}

Exemplo completo de criação de Merge Request:

\begin{terminal}
glab mr create --source-branch desenvolvimento --target-branch main \
  --title "Integra o sistema ao repositório institucional" \
  --description "Une o histórico local ao histórico do GitLab."
\end{terminal}

Todo comando tem ajuda embutida: \cmd{glab mr create --help}.

\section{Terminal ou interface web?}

Tudo o que é versionamento e merge dá para fazer no Git, localmente. A parte de colaboração no GitLab dá para fazer com o \cmd{glab}. Mesmo assim, em ambiente institucional a \textbf{interface web continua importante} para revisão visual, aprovação, auditoria, acompanhamento dos pipelines e conversa da equipe.

\begin{atencao}[As regras valem no terminal também]
O \cmd{glab} não pula regras. Se a \arq{main} exige aprovação de um responsável, \cmd{glab mr merge} vai esperar essa aprovação como a interface web esperaria.
\end{atencao}
